\chapter{Gravitational Waves}

\section*{Theory problems}

\probhead{14.1}{Gravitational Waves}{2016}{Bhubaneswar, India}{TH T11}{50}
% source id: 2016_TH_T11   tag: GW frequency at merger MS/WD/NS/BH
The first signal of gravitational waves was observed by two advanced LIGO
detectors at Hanford and Livingston, USA in September 2015. One of these
measurements (strain vs.\ time in seconds) is shown in the accompanying figure.
In this problem, we will interpret this signal in terms of a small test mass $m$
orbiting around a large mass $M$ (i.e.\ $m \ll M$), by considering several
models for the nature of the central mass.

\ioaafig{2016_TH_T11-f1.pdf}

The test mass loses energy due to the emission of gravitational waves. As a
result the orbit keeps on shrinking, until the test mass reaches the surface of
the object, or in the case of a black hole, the innermost stable circular orbit
--- ISCO --- which is given by $R_{\mathrm{ISCO}} = 3R_{\mathrm{sch}}$, where
$R_{\mathrm{sch}}$ is the Schwarzschild radius of the black hole. This is the
``epoch of merger''. At this point, the amplitude of the gravitational wave is
maximum, and so is its frequency, which is always twice the orbital frequency.
In this problem, we will only focus on the gravitational waves before the
merger, when Kepler's laws are assumed to be valid. After the merger, the form
of gravitational waves will drastically change.

\begin{description}
\item[(T11.1)] Consider the observed gravitational waves shown in the figure
  above. Estimate the time period, $T_0$, and hence calculate the frequency,
  $f_0$, of gravitational waves just before the epoch of merger. \hfill[7]
\item[(T11.2)] For any main sequence (MS) star, the radius of the star,
  $R_{\mathrm{MS}}$, and its mass, $M_{\mathrm{MS}}$, are related by a power
  law. Find the highest frequency of emitted gravitational waves,
  $f_{\mathrm{MS,max}}$, if the test mass is orbiting a main sequence star.
\item[(T11.3)] Find the highest frequency of emitted gravitational waves,
  $f_{\mathrm{WD,max}}$, if the test mass is orbiting a white dwarf. Take the
  radii of white dwarfs to be in the range $\sim 3000$ to $6000$\,km.
\item[(T11.4)] Neutron stars (NS) are a peculiar type of compact objects which
  have masses between $1$ and $3M_\odot$ and radii in the range $10$--$15$\,km.
  Find the range of frequencies of emitted gravitational waves,
  $f_{\mathrm{NS,min}}$ and $f_{\mathrm{NS,max}}$, if the test mass is orbiting
  a neutron star at a distance close to the neutron star radius.
\item[(T11.5)] If the test mass is orbiting a black hole (BH), write the
  expression for the frequency of emitted gravitational waves,
  $f_{\mathrm{BH}}$, in terms of the mass of the black hole,
  $M_{\mathrm{BH}}$, and the solar mass $M_\odot$. \hfill[5]
\item[(T11.6)] Based only on the time period (or frequency) of gravitational
  waves before the epoch of merger, determine whether the central object can be
  a main sequence star (MS), a white dwarf (WD), a neutron star (NS), or a
  black hole (BH). Estimate the mass of this object, $M_{\mathrm{obj}}$, in
  units of $M_\odot$. \hfill[8]
\end{description}

\probhead{14.2}{LIGO}{2021}{Bogot\'a (remote)}{TH Q1}{5}
% source id: 2021_TH_Q1   tag: BH-merger GW energy vs supernova
The first detection of gravitational waves GW150914 was announced in 2016 by
the collaboration LIGO (Laser Interferometer Gravitational-Wave Observatory).
The detected signal corresponds to the merger of two black holes with masses of
$35M_\odot$ and $30M_\odot$, which when joined formed a black hole of
$62M_\odot$. Ignoring the rotational energies of the black holes, you may assume
that the energy released by this process ($E_{GW}$) is emitted solely in the
form of gravitational waves, that were observed by the interferometer in 2015.
You are given that the explosion of a supernova (SN) releases
$E_{SN} = 2 \times 10^{44}$\,J.

\begin{description}
\item[(1.1)] To find out which of these two events (SN, GW) releases more
  energy, estimate the energy ratio $E_{SN}/E_{GW}$. \hfill[5.0\,pt]
\end{description}

\probhead{14.3}{LISA}{2023}{Katowice, Poland}{TH Q13}{45}
% source id: 2023_TH_Q13   tag: white-dwarf binary GW strain, chirp mass
The Laser Interferometer Space Antenna (LISA) is a proposed experiment to detect
low-frequency gravitational waves. It consists of three spacecraft arranged in
an equilateral triangle. A passing gravitational wave changes the distance
between the spacecraft, which can be precisely measured (more details are given
in the notes below).

One of the sources of low-frequency gravitational waves are compact binary star
systems, for example binary white dwarfs. Such a system was recently discovered
at a distance of $2.34$\,kpc from the Sun. The orbital period of the binary was
found to be $414.79$\,s and is changing at a rate of
$-7.49 \times 10^{-4}$\,s\,yr$^{-1}$ due to the emission of gravitational waves.

\begin{description}
\item[(a)] Check if this binary system can be detected by LISA. \hfill[25 points]
\item[(b)] Calculate the chirp mass. \hfill[5 points]
\item[(c)] Determine the masses of both components knowing that the ratio
  between the radius of one of the components to the semi-major axis of the
  orbit is $0.139$, and assuming both components follow the mass--radius
  relation for white dwarfs given in the table below. \hfill[15 points]
\end{description}

\noindent\textbf{Notes:}
\begin{enumerate}
\item A binary star system with an orbital period $P$ emits gravitational waves
  with a frequency of $f = 2/P$.
\item LISA measures a dimensionless quantity called the characteristic strain
  amplitude, $S$, given by
  \[ S = h\sqrt{f T_{\mathrm{obs}}}, \]
  where $T_{\mathrm{obs}} = 4$\,yr is the expected duration of the mission. $h$
  is the gravitational wave strain, given by
  \[ h = \frac{2(G\mathcal{M})^{5/3}(\pi f)^{2/3}}{c^4 D}, \]
  where $\mathcal{M}$ is the so-called chirp mass, $f$ is the frequency of the
  gravitational wave and $D$ is the distance to the system. If we denote the
  masses of the components by $M_1$ and $M_2$, then the chirp mass is given by
  \[ \mathcal{M} = \frac{(M_1 M_2)^{3/5}}{(M_1 + M_2)^{1/5}}. \]
  The expected sensitivity of LISA as a function of gravitational wave
  frequency is presented in the figure below.
\item The semi-major axis $a$ of the binary system changes due to the emission
  of gravitational waves at a rate
  \[ \frac{\Delta a}{\Delta t} = -\frac{64}{5}\,\frac{G^3}{c^5}\,
     \frac{M_1 M_2 (M_1 + M_2)}{a^3}. \]
\end{enumerate}

\ioaafig{2023_TH_Q13-f2.pdf}

\probhead{14.4}{Gravitational Waves}{2025}{Mumbai, India}{TH T03}{15}
% source id: 2025_TH_T03   tag: binary BH chirp, merger time
Orbiting binary black holes generate gravitational waves. Consider two black
holes, in our Galaxy, with masses $M = 36\,M_\odot$ and $m = 29\,M_\odot$,
revolving in circular orbits with orbital angular frequency $\omega$ around
their centre of mass.

\begin{description}
\item[(T03.1)] Assuming Newtonian gravity, derive an expression for the angular
  frequency, $\omega_{\mathrm{ini}}$, of the black hole orbits at a time
  $t_{\mathrm{ini}}$, when the separation between the black holes was $4.0$
  times the sum of their Schwarzschild radii, in terms of only $M$, $m$, and
  physical constants. Calculate the value of $\omega_{\mathrm{ini}}$
  (in rad\,s$^{-1}$). \hfill[5]
\item[(T03.2)] In general relativity, black holes in orbit emit gravitational
  waves with frequency $f_{\mathrm{GW}}$, such that
  $2\pi f_{\mathrm{GW}} = \omega_{\mathrm{GW}} = 2\omega$. This shrinks the
  black hole orbits, which in turn increases $f_{\mathrm{GW}}$. The rate of
  change of $f_{\mathrm{GW}}$ is
  \[ \frac{\mathrm{d}f_{\mathrm{GW}}}{\mathrm{d}t}
     = \frac{96\pi^{8/3}}{5}\,G^{5/3} c^{\beta}
       \mathcal{M}_{\mathrm{chirp}}^{\alpha/3} f_{\mathrm{GW}}^{\delta/3}, \]
  where $\mathcal{M}_{\mathrm{chirp}} = \dfrac{(mM)^{3/5}}{(m+M)^{1/5}}$ is
  called the ``chirp mass''. Find the values of $\alpha$, $\beta$ and $\delta$.
  \hfill[4]
\item[(T03.3)] Assume that the gravitational waves associated with the event
  were first detected at time $t_{\mathrm{ini}} = 0$. Derive an expression for
  the observed time of black hole merger, $t_{\mathrm{merge}}$, when
  $f_{\mathrm{GW}}$ becomes very large, in terms of $\omega_{\mathrm{ini}}$,
  $\mathcal{M}_{\mathrm{chirp}}$, and physical constants only. Calculate the
  value of $t_{\mathrm{merge}}$ (in seconds). \hfill[6]
\end{description}

\section*{Data-analysis problems}

\probhead{14.5}{Gravitational wave astronomy}{2022}{Kutaisi, Georgia}{DA 1}{45}
% source id: 2022_DA_1   tag: GW170817 chirp mass from f(t)
In this problem, we will analyze frequency vs.\ time representations of data
containing the gravitational wave event GW170817, which was the first
observation by LIGO of a merger of a binary neutron star system.

\ioaafig{2022_DA_1-f1.pdf}

\subsection*{1.1 Data from the graph}
\begin{description}
\item[(1.1.1)] You are given a semi-logarithmic graph, from which you need to
  extract values of frequency and time. Let the bottom left corner be the
  origin $(0,0)$, horizontal coordinate $x$ is measured in centimetres. Find a
  linear expression to obtain actual time $t$ from $x$. \hfill[2\,pt]
\item[(1.1.2)] Similarly, find an expression for the frequency $f$ as a
  function of the vertical coordinate $y$ measured in centimetres.
  \hfill[3\,pt]
\item[(1.1.3)] Extract 14 values of time and corresponding frequency from the
  given graph into a table. At least one of the values should correspond to a
  frequency more than $100$\,Hz. \hfill[7\,pt]
\end{description}

\subsection*{1.2 Calculate system parameters}
The most plausible explanation for this evolution of frequency is the
in-spiraling of two orbiting masses, $m_1$ and $m_2$, due to gravitational-wave
emission. At the lower frequencies, such evolution is characterized by the
``chirp mass''
\[ M_{\mathrm{chirp}} = \frac{(m_1 m_2)^{3/5}}{(m_1+m_2)^{1/5}}
   = \frac{c^3}{G}\left[\frac{5}{96}\,\pi^{-8/3} f^{-11/3}\dot{f}\right]^{3/5}, \]
where $f$ and $\dot{f}$ are the observed frequency and its rate of change with
respect to time, and $G$ and $c$ are the gravitational constant and the speed
of light.

\begin{description}
\item[(1.2.1)] From the equation above, derive an expression for the frequency
  dependence on time. \hfill[3\,pt]\\
  \textit{Note:} if $x^n \dot{x} = k$, then $\dfrac{x^{n+1}}{n+1} = kt + C$,
  where $k$, $n$ and $C$ are constants.
\item[(1.2.2)] By plotting the appropriate graph, find the chirp mass and
  estimate its error caused by uncertainty in the slope calculation.
  \hfill[22\,pt]
\end{description}

We realize that what is measured by the ground-based GW detectors are actually
the detector-frame masses, which are related to the source-frame masses by
\[ m_{\mathrm{detector}} = (1+z)\,m, \]
where $z$ is the redshift of the binary.

\begin{description}
\item[(1.2.3)] If the host galaxy NGC 4993 has a redshift $z = 0.009783$, find
  the source-frame chirp mass. \hfill[2\,pt]
\item[(1.2.4)] The mass ratio $q = m_1/m_2$ is much harder to measure. Advanced
  waveform analysis constrains this ratio; using the chirp mass found above,
  determine the range of possible values of the individual masses $m_1$ and
  $m_2$. \hfill[6\,pt]
\end{description}

\clearpage
\section*{Solutions to Chapter 14}

\solhead{14.1}{Gravitational Waves}
% source id: 2016_TH_T11
\begin{description}
\item[(T11.1)] From the graph, just before the peak of emission, the time
  period of gravitational waves is approximately $(0.007 \pm 0.004)$\,s.
  \[ T_0 \approx 0.007\,\mathrm{s} \]
  \hfill(2 points)

  That is, the frequency of the gravitational waves is
  $f_0 \approx 142.86$\,Hz. \hfill(1 point)

\item[(T11.2)] Since $m \ll M$ then, by Kepler's third law
  \[ f_{\mathrm{orbital}} = \frac{1}{2\pi}\sqrt{\frac{GM}{r^3}} \]
  \hfill(4 points)

  Hence the frequency of the gravitational waves is
  \[ f_{\mathrm{grav}} = 2 f_{\mathrm{orbital}}
     = \frac{1}{\pi}\sqrt{\frac{GM}{r^3}} \]
  \hfill(1 point)

  The frequency will be maximum when $r = R_{\mathrm{MS}}$. \hfill(1 point)

  For main sequence stars,
  \[ \frac{R_{\mathrm{MS}}}{R_\odot}
     = \left(\frac{M_{\mathrm{MS}}}{M_\odot}\right)^{\alpha}
     \qquad\therefore\quad R_{\mathrm{MS}}
     = R_\odot \left(\frac{M_{\mathrm{MS}}}{M_\odot}\right)^{\alpha} \]
  \hfill(1 point)
  \begin{align*}
    \therefore f_{\mathrm{MS}}
      &= \frac{1}{\pi}\sqrt{\frac{G M_{\mathrm{MS}}}{R_\odot^3}}
         \left(\frac{M_\odot}{M_{\mathrm{MS}}}\right)^{3\alpha/2} \\
      &= \frac{1}{\pi}\sqrt{\frac{G M_\odot}{R_\odot^3}}
         \left(\frac{M_\odot}{M_{\mathrm{MS}}}\right)^{(3\alpha-1)/2} \\
    f_{\mathrm{MS}}
      &= \frac{1}{\pi}\sqrt{\frac{G M_\odot}{R_\odot^3}}
         \left(\frac{M_{\mathrm{MS}}}{M_\odot}\right)^{(1-3\alpha)/2}
  \end{align*}
  \hfill(3 points)

\item[(T11.3)] For possible values of $\alpha$ given in the question, the
  exponent $\frac{1-3\alpha}{2}$ is negative. Thus, if
  $M_{\mathrm{MS}} > M_\odot$, the frequency will be smaller. Thus, for highest
  possible frequency coming from a main sequence star, you should take lowest
  possible mass i.e.\ $\alpha = 1.0$. \hfill(4 points)
  \begin{align*}
    f_{\mathrm{MS,max}}
      &= \frac{1}{\pi}\sqrt{\frac{G M_\odot}{R_\odot^3}}
         \left(\frac{M_{\mathrm{MS}}}{M_\odot}\right)^{\left(\frac{1-3\times 1}{2}\right)} \\
      &= \frac{1}{\pi}\sqrt{\frac{G M_\odot}{R_\odot^3}}
         \times \frac{M_\odot}{M_{\mathrm{MS}}}
  \end{align*}
  \hfill(2 points)

  The frequency of gravitational waves will be given by,
  \begin{align*}
    f_{\mathrm{MS,max}}
      &= \frac{1}{\pi}\sqrt{\frac{6.674 \times 10^{-11} \times 1.989 \times 10^{30}}
         {(6.955 \times 10^{8})^{3}}} \times \frac{1}{0.08} \\
    f_{\mathrm{MS,max}} &= 2.5\,\mathrm{mHz}
  \end{align*}
  \hfill(3 points)

\item[(T11.4)] The maximum frequency would be when $r = R_{\mathrm{WD}}$.
  \hfill(1 point)

  We use the notation $R_{WD\odot}$ for the radius of a solar mass white
  dwarf. Then for white dwarfs
  % sic: the official solution prints the mass ratio below as M_WD/M_sun;
  % from R \propto M^{-1/3} it should be M_sun/M_WD, which is what the
  % subsequent lines (and the numerical answer) actually use.
  \[ R_{WD}^3 = R_{WD\odot}^3 \frac{M_{WD}}{M_\odot} \]
  \hfill(1 point)
  \begin{align*}
    f_{\mathrm{WD}}
      &= \frac{1}{\pi}\sqrt{\frac{G M_{\mathrm{WD}}}{R_{WD}^3}} \\
      &= \frac{1}{\pi}\sqrt{\frac{G M_\odot}{R_{WD\odot}^3}}\,
         \frac{M_{WD}}{M_\odot}
  \end{align*}
  \hfill(2 points)

  For maximum frequency, we have to take highest white dwarf mass.
  \hfill(2 points)
  \begin{align*}
    f_{\mathrm{WD,max}}
      &= 2.600 \times 10^{-6} \times
         \sqrt{\frac{1.989 \times 10^{30}}{(6000 \times 10^{3})^{3}}}
         \times 1.44 \\
      &= 2.600 \times 10^{-6} \times 95.96 \times 10^{3} \times 1.44 \\
    f_{\mathrm{WD,max}} &= 0.359\,\mathrm{Hz}
  \end{align*}
  \hfill(2 points)

\item[(T11.5)]
  \[ f_{grav} = \frac{1}{\pi}\sqrt{\frac{G M_{\mathrm{NS}}}{R_{\mathrm{NS}}^3}} \]
  The lowest possible frequency is when $M_{\mathrm{NS}}$ is lowest and
  $R_{\mathrm{NS}}$ is the highest. \hfill(2 points)

  For $M_{\mathrm{NS}} = M_\odot$ and $R = 15$\,km, we get
  \[ f_{\mathrm{NS,min}} = 1.996\,\mathrm{kHz} \]
  \hfill(2 points)

  Similarly, the largest possible frequency is when $M_{\mathrm{NS}}$ is the
  largest and $R_{\mathrm{NS}}$ is the smallest. \hfill(2 points)

  For $M_{\mathrm{NS}} = 3M_\odot$ and $R = 10$\,km, we get
  \[ f_{\mathrm{NS,max}} = 6.352\,\mathrm{kHz} \]
  \hfill(2 points)

\item[(T11.6)] For black holes, we have to consider $R_{ISCO}$.
  \hfill(2 points)

  Hence the equation will be,
  \[ f_{BH} = \frac{1}{\pi} \times
     \sqrt{\frac{G M_\odot}{27 R_{sch-\odot}^3}} \times
     \frac{M_\odot}{M_{BH}} \]
  \hfill(2 points)
  \[ f_{BH} = 4.396\,\mathrm{kHz} \times \frac{M_\odot}{M_{BH}} \]
  \hfill(3 points)

\item[(T11.7)] We found the frequency of the LIGO-detected wave to be
  166.67\,Hz just before merger.
  % sic: the official solution says 166.67 Hz here, although part (T11.1)
  % found 142.86 Hz (and 142.86 is used in the mass estimate below).
  As per our analysis above, only black holes can lead to emission in this
  frequency range: Black Hole. \hfill(2 points)

  By using corresponding expression,
  \[ M_{\mathrm{obj}} = \frac{4396}{142.86} M_\odot \approx 31 M_\odot \]
  \hfill(3 points)
\end{description}

\solhead{14.2}{LIGO}
% source id: 2021_TH_Q1
\begin{description}
\item[(1.1)] Based on Einstein's relationship, $E = mc^2$. Doing a mass-energy
  balance:
  % note: the source writes the speed of light as a capital C throughout.
  \[ 35 M_\odot c^2 + 30 M_\odot c^2 = 62 M_\odot c^2 + E_{GW} \]
  \hfill(1 point)
  \[ 3 M_\odot c^2 = E_{GW} \]
  \hfill(1 point)
  \[ M_\odot = 1.988 \times 10^{30}\,\mathrm{kg} \]
  \[ E_{GW} = 3\,(1.988 \times 10^{30}\,\mathrm{kg})
     (9 \times 10^{16}\,\mathrm{m^2\,s^{-2}}) \]
  \hfill(1 point)
  \[ E_{GW} = 5.36 \times 10^{47}\,\mathrm{J} \]
  \hfill(1 point)

  Now using $E_{SN} = 2 \times 10^{44}$\,J we have the ratio:
  \[ \frac{E_{SN}}{E_{GW}} = 3.7 \times 10^{-4} \]
  \hfill(1 point)

  GW150914 released approximately 2865 times more energy than a supernova
  explosion. $E_{SN} \ll E_{GW}$.
\end{description}

\solhead{14.3}{LISA}
% source id: 2023_TH_Q13
% Note: the official solution is divided into Part (a) and Part (b) only;
% its Part (b) covers both parts (b) and (c) of the problem statement.
\begin{description}
\item[Part (a)] To determine whether the system can be detected by LISA, we
  need to determine two quantities: the gravitational-wave frequency and the
  characteristic strain amplitude.

  It is straightforward to calculate the gravitational-wave frequency:
  \[ f = \frac{2}{P} = \frac{2}{414.79} = 4.8 \times 10^{-3}\,\mathrm{Hz}. \]
  This frequency is within the LISA band and close to the maximum LISA
  sensitivity. \hfill(2 points)

  To estimate the characteristic strain amplitude, we need to know the chirp
  mass $\mathcal{M}$. The other required quantities (such as the gravitational
  wave frequency, distance, and duration of the LISA observations) are already
  known.

  The orbital period change rate $\Delta P/\Delta t$ is given in the problem.
  We need to link it to $\Delta a/\Delta t$ which we know from Note 3 is
  linked to the mass function. From Kepler's third law, we know that:
  \[ \frac{a^3}{P^2} = \frac{G(M_1 + M_2)}{4\pi^2}, \]
  where $M_1$ and $M_2$ are masses of both components of the system. If, due
  to the emission of gravitational waves, the semi-major axis changes from $a$
  to $a + \Delta a$, then the orbital period changes from $P$ to
  $P + \Delta P$, as the masses of both components are constant. Then:
  \[ \frac{a^3}{P^2} = \frac{(a+\Delta a)^3}{(P+\Delta P)^2}
     = \frac{a^3(1+\Delta a/a)^3}{P^2(1+\Delta P/P)^2}
     = \frac{a^3}{P^2}\left(1 + \frac{3\Delta a}{a}
       - \frac{2\Delta P}{P}\right), \]
  hence:
  \[ \frac{\Delta P}{P} = \frac{3}{2}\,\frac{\Delta a}{a}. \]
  Here, we used the fact that $(1+x)^n \approx 1 + nx$ for $x \ll 1$.

  Therefore:
  \begin{align*}
    \frac{\Delta P}{\Delta t}
      &= \frac{3P}{2a}\frac{\Delta a}{\Delta t}
       = -\frac{3}{2}\cdot\frac{64}{5}\,\frac{P}{a}\,\frac{G^3}{c^5}\,
         \frac{M_1 M_2 (M_1+M_2)}{a^3} \\
      &= -\frac{96}{5}\,\frac{G^3}{c^5}\,\frac{P}{a}\,
         \frac{M_1 M_2 (M_1+M_2)}{G(M_1+M_2)P^2}\cdot 4\pi^2 \\
      &= -\frac{96}{5}(2\pi)^2\,\frac{G^2}{c^5}\,\frac{M_1 M_2}{aP}
       = -\frac{96}{5}(2\pi)^2\,\frac{G^2}{c^5}\,\frac{M_1 M_2}{P}\,
         \frac{(2\pi)^{2/3}}{G^{1/3}(M_1+M_2)^{1/3}P^{2/3}} \\
      &= -\frac{96}{5}(2\pi)^{8/3}\,\frac{G^{5/3}}{c^5}\,
         \frac{M_1 M_2}{(M_1+M_2)^{1/3}}\,\frac{1}{P^{5/3}} \\
      % note: the source prints the expression on the line above twice.
      &= -\frac{96}{5c^5}(2\pi)^{8/3}
         \left(\frac{G\mathcal{M}}{P}\right)^{5/3}
       = -\frac{192\pi}{5c^5}(G\mathcal{M})^{5/3}
         \left(\frac{P}{2\pi}\right)^{-5/3},
  \end{align*}
  Thus, by knowing the orbital period and its rate of change from
  observations, we can determine the chirp mass:
  \[ \mathcal{M} = \left(\frac{5}{192\pi}\right)^{3/5}\frac{c^3}{G}\,
     \frac{P}{2\pi}\left(-\frac{\Delta P}{\Delta t}\right)^{3/5}
     = 0.319\,M_\odot, \]
  and find the characteristic strain amplitude $h$.

  Alternatively, if we notice that the characteristic strain amplitude $h$
  depends on $(G\mathcal{M})^{5/3}$ which we can get directly from the rate of
  change of the period:
  \[ (G\mathcal{M})^{5/3} = -\frac{\Delta P}{\Delta t}
     \left(\frac{P}{2\pi}\right)^{5/3}\frac{5c^5}{192\pi}, \]
  we can skip calculating the chirp mass and get the characteristic strain
  amplitude directly, which is all we need to check if the binary can be
  detected.

  Either way, the gravitational wave strain is:
  \[ h = -\frac{5c}{192\pi^2}\,P\,\frac{\Delta P}{\Delta t}\,\frac{1}{D}. \]
  If we plug in the numerical values, we get $h = 1.1 \times 10^{-22}$ and
  $S = 8.4 \times 10^{-20}$. Checking the plot, this is above the expected
  sensitivity of LISA at 5\,mHz. Thus, this object should be detected by LISA.

\item[Part (b)] To determine the masses of both components $M_1$ and $M_2$ we
  need two simultaneous $M_1 - M_2$ relations. The expression for the chirp
  mass:
  \[ \mathcal{M} = \frac{(M_1 M_2)^{3/5}}{(M_1 + M_2)^{1/5}}, \]
  gives us one relation, and we can calculate the chirp mass of the observed
  system from the rate of change of the period, if that was not already done
  in part (a):
  \[ \mathcal{M} = \left(\frac{5}{192\pi}\right)^{3/5}\frac{c^3}{G}\,
     \frac{P}{2\pi}\left(-\frac{\Delta P}{\Delta t}\right)^{3/5}
     = 0.319\,M_\odot. \]
  The second relation can be obtained from the mass--radius relation for white
  dwarfs given in the table and Kepler's third law.

  We are told that for one of the components (call it `1'),
  \[ a = \frac{R_1}{0.139}. \]
  From Kepler's third law we know that:
  \[ M_1 + M_2 = \left(\frac{a}{1\,\mathrm{au}}\right)^3
     \left(\frac{P}{1\,\mathrm{yr}}\right)^2, \]
  which will give us the mass of the second component $M_2$ from the mass of
  the first, $M_1$.

  From here, it is not possible to derive an analytical formula for the masses
  of the components. Instead, we need to use numerical methods to estimate the
  result.

  Taking the masses and radii listed in the given table as $M_1$ and $R_1$, we
  can obtain $a$ and thus $M_1 + M_2$, $M_2$ and finally $\mathcal{M}$ for
  each mass:

  \begin{center}
  \begin{tabular}{cccccc}
    \hline
    $M_1$ ($M_\odot$) & $R$ ($R_\odot$) & $a$ ($R_\odot$) & $M_1+M_2$ &
    $M_2$ & $\mathcal{M}$ \\
    \hline
    0.48 & 0.0144 & 0.104 & 0.647 & 0.167 & 0.240 \\
    0.50 & 0.0147 & 0.106 & 0.689 & 0.189 & 0.261 \\
    0.52 & 0.0150 & 0.108 & 0.732 & 0.212 & 0.283 \\
    0.54 & 0.0153 & 0.110 & 0.776 & 0.236 & 0.306 \\
    0.56 & 0.0156 & 0.112 & 0.823 & 0.263 & 0.329 \\
    0.58 & 0.0159 & 0.114 & 0.871 & 0.291 & 0.354 \\
    0.60 & 0.0162 & 0.117 & 0.922 & 0.322 & 0.379 \\
    0.62 & 0.0165 & 0.119 & 0.974 & 0.354 & 0.405 \\
    0.64 & 0.0168 & 0.121 & 1.028 & 0.388 & 0.431 \\
    \hline
  \end{tabular}
  \end{center}

  The actual chirp mass is $\mathcal{M} = 0.319\,M_\odot$. Therefore, by
  linear interpolation or graphically, we estimate $M_1 = 0.55\,M_\odot$ and
  $M_2 = 0.25\,M_\odot$.

  (The student does not need to calculate $\mathcal{M}$ for all values of
  $M_1$.)

  \ioaafig{2023_TH_Q13-s1.pdf}
\end{description}

\solhead{14.4}{Gravitational Waves}
% source id: 2025_TH_T03
\begin{description}
\item[(T03.1)] Assuming Newtonian gravity, for black hole masses $M$ and $m$
  at distances $R$ and $r$, respectively, from their centre of mass,
  \begin{align*}
    \omega^2 r &= \frac{GM}{(R+r)^2} \\
    \omega &= \left[\frac{GM}{r(R+r)^2}\right]^{1/2}
  \end{align*}
  \hfill(1 point)

  Since $MR = mr$, we can rewrite this as
  \[ \omega = \left[\frac{G(M+m)}{(R+r)^3}\right]^{1/2} \]
  \hfill(1 point)

  Let $R_{\mathrm{s}}$ and $r_{\mathrm{s}}$ be the respective Schwarzschild
  radii of the black holes of mass $M$ and $m$. Expressing the separation
  $(r+R) = 4.0 \times (r_{\mathrm{s}} + R_{\mathrm{s}})$ at
  $t_{\mathrm{ini}}$, we obtain,
  \begin{align*}
    \omega_{\mathrm{ini}}
      &= \left[\frac{G(M+m)}{4^3 (R_{\mathrm{s}}+r_{\mathrm{s}})^3}\right]^{1/2}
       = \left[\frac{c^6 G(M+m)}{8 \times 4^3 G^3 (M+m)^3}\right]^{1/2} \\
      &= \frac{c^3}{16\sqrt{2}\,G(M+m)} \\
    \omega_{\mathrm{ini}} &= \frac{c^3}{16\sqrt{2}\,G(M+m)}
  \end{align*}
  \hfill(2 points)

  Substituting the values of the masses,
  \[ \omega_{\mathrm{ini}} = 1.4 \times 10^2\,\mathrm{rad\,s^{-1}}. \]
  \hfill(1 point)

\item[(T03.2)] We use dimensional analysis to calculate the exponents in the
  above equation. The dimension of $\mathcal{M}_{\mathrm{chirp}}$ is [M].

  The dimensions of the left and right hand sides of the equation are
  \[ [\mathrm{T}]^{-2} = [\mathrm{T}]^{-\delta/3}
     \left([\mathrm{L}]^3[\mathrm{T}]^{-2}[\mathrm{M}]^{-1}\right)^{5/3}
     \left([\mathrm{L}][\mathrm{T}]^{-1}\right)^{\beta}
     [\mathrm{M}]^{\alpha/3} \]
  \hfill(1 point)

  Equating the dimensions for T, L, and M, we get
  \begin{align*}
    -2 &= -\frac{\delta}{3} - \frac{10}{3} - \beta \\
    0 &= 5 + \beta \\
    0 &= \frac{\alpha}{3} - \frac{5}{3},
  \end{align*}
  respectively. Solving, we get $\alpha = 5$, $\beta = -5$ and $\delta = 11$.
  \hfill(3 points)

\item[(T03.3)] From (T03.2), we get
  \[ \frac{\mathrm{d}f_{\mathrm{GW}}}{\mathrm{d}t}
     = \frac{96\pi^{8/3}}{5}\,\frac{G^{5/3}}{c^5}\,
       \mathcal{M}_{\mathrm{chirp}}^{5/3} f_{\mathrm{GW}}^{11/3}. \]
  \hfill(1 point)

  Integrating the equation for $f_{\mathrm{GW}}$, we obtain
  \[ \int_{\omega_{\mathrm{ini}}/\pi}^{\infty}
     f_{\mathrm{GW}}^{-11/3}\,\mathrm{d}f_{\mathrm{GW}}
     = \int_{0}^{t_{\mathrm{merge}}} \frac{96\pi^{8/3}}{5}\,
       \frac{G^{5/3}}{c^5}\,\mathcal{M}_{\mathrm{chirp}}^{5/3}\,\mathrm{d}t \]
  \hfill(1 point)
  \[ \left.-\frac{3}{8} f_{\mathrm{GW}}^{-8/3}
     \right|_{\omega_{\mathrm{ini}}/\pi}^{\infty}
     = \frac{96\pi^{8/3}}{5}\,\frac{G^{5/3}}{c^5}\,
       \mathcal{M}_{\mathrm{chirp}}^{5/3}\,t_{\mathrm{merge}} \]
  \hfill(1 point)
  \begin{align*}
    t_{\mathrm{merge}}
      &= \frac{3}{8}\left[\frac{\omega_{\mathrm{ini}}}{\pi}\right]^{-8/3}
         \frac{5}{96\pi^{8/3}}\,\frac{c^5}{G^{5/3}}\,
         \frac{1}{\mathcal{M}_{\mathrm{chirp}}^{5/3}} \\
      &= \frac{5}{256}\,\omega_{\mathrm{ini}}^{-8/3}\,
         \frac{c^5}{(G\mathcal{M}_{\mathrm{chirp}})^{5/3}}. \\
    t_{\mathrm{merge}}
      &= \frac{5}{256}\,\omega_{\mathrm{ini}}^{-8/3}\,
         \frac{c^5}{(G\mathcal{M}_{\mathrm{chirp}})^{5/3}}
  \end{align*}
  \hfill(1 point)

  We calculate
  $\mathcal{M}_{\mathrm{chirp}} = \dfrac{(mM)^{3/5}}{(m+M)^{1/5}}
   = 28\,\mathrm{M}_\odot$. \hfill(1 point)

  Substituting the values of $\mathcal{M}_{\mathrm{chirp}}$, and
  $\omega_{\mathrm{ini}}$ we get the value of
  $t_{\mathrm{merge}} = 0.10$\,s. \hfill(1 point)
\end{description}

\solhead{14.5}{Gravitational wave astronomy}
% source id: 2022_DA_1
% Note: part labels below are those of the official solution ((1.1)-(3.5)).
% The statement transcribed in the book is the released exam version, whose
% numbering (1.1.1, 1.2.1, ...) differs slightly and which omits the
% distance part (2.4) and the whole section 3 (speed of the gravitational
% wave); the full official solution is transcribed here.

\noindent\textbf{1\quad Data from the graph}
\begin{description}
\item[(1.1)] To make the scaling as accurate as possible, we take the time
  values far apart. When $x = 0.0$\,cm $t = -30$\,s, using the ruler we can
  also measure that at $t = 0$\,s we have $x = 11.1$\,cm. \hfill(1 point)

  This leads to
  \begin{align*}
    0.0\,\mathrm{cm} &= (-30\,\mathrm{s})A + B \\
    11.1\,\mathrm{cm} &= (0\,\mathrm{s})A + B = B \\
    A &= \frac{11.1\,\mathrm{cm}}{30\,\mathrm{s}} = 0.37\,\mathrm{cm/s} \\
    \therefore t &= \frac{x - 11.1}{0.37}\,\mathrm{s}
  \end{align*}
  \hfill(1 point)

\item[(1.2)] Analogues to the previous scaling,
  \begin{align*}
    0.0\,\mathrm{cm} &= C\log(30\,\mathrm{Hz}) + D \\
    5.3\,\mathrm{cm} &= C\log(500\,\mathrm{Hz}) + D
  \end{align*}
  \hfill(1 point)
  \begin{align*}
    5.3\,\mathrm{cm} &= C\log\left(\frac{500}{30}\right) \\
    C &= 4.34\,\mathrm{cm} \\
    D &= -6.41\,\mathrm{cm} \\
    \therefore f &= 10^{\left(\frac{y+6.41}{4.34}\right)}\,\mathrm{Hz}
  \end{align*}
  \hfill(2 points)

\item[(1.3)] It is not necessary to use $X$-scaling as students can choose
  convenient points from the graph.

  \begin{center}
  \begin{tabular}{ccccc}
    \hline
    $x$ (cm) & $y$ (cm) & $t$ (s) & $f$ (Hz) & $f^{-8/3}$ ($\mathrm{Hz}^{-8/3}$) \\
    \hline
    - & 0.5 & $-28.0$ & 39.1 & 5.67E-05 \\
    - & 0.5 & $-26.0$ & 39.1 & 5.67E-05 \\
    - & 0.6 & $-24.0$ & 41.3 & 4.92E-05 \\
    - & 0.6 & $-22.0$ & 41.3 & 4.92E-05 \\
    - & 0.7 & $-20.0$ & 43.5 & 4.27E-05 \\
    - & 0.8 & $-18.0$ & 45.9 & 3.71E-05 \\
    - & 0.9 & $-16.0$ & 48.4 & 3.22E-05 \\
    - & 1.0 & $-14.0$ & 51.0 & 2.79E-05 \\
    - & 1.1 & $-12.0$ & 53.8 & 2.43E-05 \\
    - & 1.2 & $-10.0$ & 56.7 & 2.11E-05 \\
    - & 1.4 & $-8.0$  & 63.1 & 1.59E-05 \\
    - & 1.6 & $-6.0$  & 70.1 & 1.20E-05 \\
    - & 1.8 & $-4.0$  & 78.0 & 0.90E-05 \\
    - & 2.3 & $-2.0$  & 101.7 & 0.43E-05 \\
    - & 3.4 & 0.0     & 182.4 & 0.09E-05 \\
    \hline
  \end{tabular}
  \end{center}
\end{description}

\noindent\textbf{2\quad Calculate system parameters}
\begin{description}
\item[(2.1)]
  \begin{align*}
    \frac{96}{5}\pi^{8/3}
      \left(\frac{G M_{chirp}}{c^3}\right)^{5/3} &= f^{-11/3}\dot{f} \\
    \therefore \left[\frac{96}{5}\pi^{8/3}
      \left(\frac{G M_{chirp}}{c^3}\right)^{5/3}\right]t
      + \left[\frac{3}{8}\right] f^{-8/3} + C &= 0
  \end{align*}
  \hfill(3 points)

\item[(2.2)]
  \ioaafig{2022_DA_1-s1.pdf}
  Calculation of $f^{-8/3}$ values \hfill(3 points)

  Properly drawn graph \hfill(6 points)
  \begin{align*}
    \mathrm{slope} &= \frac{8}{3}\left[\frac{96}{5}\pi^{8/3}
      \left(\frac{G M_{chirp}}{c^3}\right)^{5/3}\right] \\
    M_{chirp} &= \left[\left(\frac{5}{256}\right)^{3/5}
      \frac{c^3}{G}\,\pi^{-8/5}\right](\mathrm{slope})^{3/5}
  \end{align*}
  \hfill(2 points)
  \[ \mathrm{slope} = 2.10 \times 10^{-6} \]
  \hfill(2 points)
  \[ \therefore M_{chirp} = 2.39 \times 10^{30}\,\mathrm{kg} = 1.20 M_\odot \]
  \hfill(2 points)

  with 95\% confidence interval
  \begin{align*}
    \Delta\mathrm{slope} &= 9 \times 10^{-8} \\
    \frac{\Delta M_{chirp}}{M_{chirp}}
      &= \frac{3}{5}\,\frac{\Delta\mathrm{slope}}{\mathrm{slope}} = 0.04 \\
    \Delta M_{chirp} &= 0.03 M_\odot
  \end{align*}

\item[(2.3)]
  \[ M_{chirp} = \frac{M_{chirp\,\mathrm{detector}}}{1+z}
     = \frac{1.20 M_\odot}{1.009783} = 1.19 M_\odot \]
  \hfill(2 points)

\item[(2.4)] According to the Hubble law,
  \[ D = \frac{zc}{H} = \frac{0.009783 \times 3 \times 10^5}{70}
     = 41.9\,\mathrm{Mpc} \]
  \hfill(2 points)

\item[(2.5)]
  \begin{align*}
    M_{chirp} &= \frac{(m_1 m_2)^{3/5}}{(m_1+m_2)^{1/5}} \\
    \therefore M_{chirp} &= m_2 \sqrt[5]{\frac{q^3}{1+q}}
  \end{align*}
  \hfill(2 points)

  This is an increasing function of $q$
  \begin{align*}
    \therefore m_{2,min} &= M_{chirp} 2^{1/5} = 1.37 M_\odot \\
    m_{2,max} &= 1.60 M_\odot
  \end{align*}
  \hfill(1 point)

  Also,
  \[ M_{chirp} = \frac{m_1}{q}\sqrt[5]{\frac{q^3}{1+q}}
     = \frac{m_1}{\sqrt[5]{q^2(1+q)}} \]
  \hfill(2 points)
  \begin{align*}
    \therefore m_{1,max} &= M_{chirp} 2^{1/5} = 1.37 M_\odot \\
    m_{2,min} &= 1.17 M_\odot
    % sic: the official solution labels this second value m_{2,min};
    % it is the minimum of the primary mass, m_{1,min}.
  \end{align*}
  \hfill(1 point)
\end{description}

\noindent\textbf{3\quad Speed of the gravitational wave}
\begin{description}
\item[(3.1)]
  \ioaafig{2022_DA_1-s2.pdf}

\item[(3.2)] It appears that the event is starting from $\Delta t = 1.8$\,s.
  Some students may also write $\Delta t = 2.2$\,s as that is when the signal
  is truly above fluctuations seen before the event. We will accept both the
  answers. \hfill(1 point)

\item[(3.3)] By simple substitution we can calculate that
  \[ \frac{\Delta v}{v_{EM}} = \frac{\Delta t\, v_{GW}}{D}
     \approx \frac{\Delta t\, c}{D} \]
  \hfill(1 point)

\item[(3.4)]
  \[ \frac{\Delta v}{v_{EM}} = 4.2 \times 10^{-16}
     \text{ to } 5.1 \times 10^{-16} \]
  \hfill(1 point)

\item[(3.5)]
  \[ \Delta t' = \Delta t - 10 = -8.2\,\mathrm{s} \text{ to } 7.8\,\mathrm{s} \]
  % sic: the official solution omits the minus sign; the upper value should
  % be -7.8 s (= 2.2 s - 10 s).
  From this time delay
  \[ \frac{\Delta v}{v_{EM}} = -1.9 \times 10^{-15}
     \text{ to } 1.8 \times 10^{-15} \]
  % sic: correspondingly, the upper value should be -1.8e-15.
  \hfill(1 point)
\end{description}
